\chapter{IMPLEMENTATION}

\section{Introduction}

This chapter describes the implementation details of the proposed 
Machine Learning-based message prioritization and congestion-aware 
rate control system for UAV communication networks. The implementation 
focuses on developing a simulation-based environment to evaluate the 
performance of the proposed communication model.

The system is implemented using software-based tools that simulate 
message transmission, congestion monitoring, and priority-based 
scheduling in UAV communication systems. The implementation includes 
message generation, feature extraction, machine learning model 
training, congestion monitoring, adaptive rate control, and message 
scheduling.

The purpose of the implementation phase is to convert the designed 
methodology into a functional system that can be tested and analyzed 
under different network conditions.

%-------------------------------------------------------------

\section{Implementation Environment}

The proposed system is implemented using a software-based simulation 
environment that supports Machine Learning and network modeling. 
The tools and technologies used during implementation are selected 
based on their efficiency, reliability, and compatibility with 
communication system modeling.

The implementation environment consists of the following components:

\begin{itemize}

\item Programming Language: Python

\item Development Platform: Python IDE

\item Machine Learning Library: Scikit-learn

\item Data Processing Library: NumPy and Pandas

\item Visualization Library: Matplotlib


\end{itemize}

Python is used as the primary programming language due to its 
simplicity and strong support for Machine Learning libraries. 
Scikit-learn is used for implementing the Random Forest classifier 
used in message prioritization.

%-------------------------------------------------------------

\section{Dataset Preparation}

Dataset preparation is an important step in the implementation of 
the Machine Learning-based classification system. The dataset 
contains MAVLink message attributes that represent communication 
behavior in UAV networks.

Each dataset entry represents a message and includes several 
features that influence its priority level. The dataset is synthetic and is 
prepared by combining multiple datasets from kaggle.

The major features included in the dataset are:

\begin{itemize}

\item Message Type

\item Payload Size

\item Transmission Frequency

\item System Status

\item Mission Phase

\item Emergency Flag

\end{itemize}

The dataset is divided into training and testing sets to evaluate 
the performance of the Machine Learning model. Typically, 
approximately 75\% of the data is used for training, and 
25\% is used for testing.

%-------------------------------------------------------------

\section{Machine Learning Model Implementation}

The Machine Learning model used in this project is based on the 
Random Forest classification algorithm. The model is implemented 
using the Scikit-learn library.

The implementation process includes the following steps:

\begin{itemize}
\item Data preprocessing and cleaning
\item Feature selection
\item Training the Random Forest classifier
\item Testing the trained model
\item Evaluating classification performance
\end{itemize}

Data preprocessing involves removing incomplete entries and 
normalizing feature values. Feature selection ensures that only 
relevant attributes are used during classification.

Each message $m_i \in \mathcal{M}$ is represented using a 
feature vector that captures its key operational attributes:

\begin{equation}
\mathbf{x}_i = [m_t,\ s,\ p_s,\ f_r,\ c,\ b,\ m_p,\ e]
\label{eq:feature_vector}
\end{equation}

where $m_t$ represents the message type, $s$ denotes the source 
node, $p_s$ is the payload size, $f_r$ is the transmission 
frequency, $c$ represents the current congestion level, $b$ 
denotes the battery state, $m_p$ indicates the mission phase, 
and $e$ is a binary indicator representing emergency conditions.

A machine learning model is employed to classify each message 
into a priority class using the mapping function:

\begin{equation}
\pi_i = f(\mathbf{x}_i), \quad \pi_i \in \{H, M, L\}
\label{eq:priority_mapping}
\end{equation}

where $H$, $M$, and $L$ correspond to high, medium, and low 
priority levels, respectively. This classification introduces 
semantic awareness into the communication process, enabling 
differentiated handling of messages based on their operational 
significance.

\begin{algorithm}[H]
\caption{Machine Learning-Based Message Priority Classification}
\begin{algorithmic}[1]
\State \textbf{Input:} MAVLink message stream $\mathcal{M}$
\State \textbf{Output:} Priority label $\pi_i$ for each message

\For{each message $m_i \in \mathcal{M}$}
    \State Extract feature vector $\mathbf{x}_i$ (Eq.~(\ref{eq:feature_vector}))
    \State Predict priority $\pi_i \leftarrow f(\mathbf{x}_i)$ (Eq.~(\ref{eq:priority_mapping}))
    \State Assign $\pi_i \in \{H, M, L\}$
\EndFor

\end{algorithmic}
\end{algorithm}

After preprocessing, the Random Forest model is trained using 
the training dataset. The trained model is then tested using 
the testing dataset to evaluate classification accuracy.

Performance metrics such as accuracy, precision, recall, and 
F1-score are calculated to measure model effectiveness 
\cite{breiman2001randomforest}.

For the combined scheduling and aggregation model, the trained policy 
is not used in isolation. Instead, it is embedded inside a closed-loop 
execution pipeline that repeatedly exchanges state and action between 
the policy and the simulation engine. This is the implementation basis 
for the reinforcement learning extension over the static scheduling 
baseline.

%-------------------------------------------------------------






%-------------------------------------------------------------
\section{Congestion Monitoring Implementation}

The congestion monitoring module is implemented to observe network 
conditions during message transmission. This module continuously 
monitors queue length, packet loss rate, and bandwidth usage.

The congestion level is quantified using a normalized congestion 
factor $c \in [0, 1]$, defined as:

\begin{equation}
c = \frac{Q_{current}}{Q_{max}}
\label{eq:congestion_factor}
\end{equation}

\noindent where $Q_{current}$ denotes the current queue occupancy 
and $Q_{max}$ represents the maximum queue capacity. When $c$ 
exceeds the predefined congestion threshold $c_{th}$, the system 
identifies network congestion and triggers adaptive rate control 
actions.

The congestion estimation logic is implemented using conditional 
checks and threshold-based monitoring techniques. The congestion 
state $\mathcal{C}$ is determined as:

\begin{equation}
\mathcal{C} = 
\begin{cases}
\text{Congested} & \text{if } c > c_{th} \\
\text{Normal}    & \text{if } c \leq c_{th}
\end{cases}
\label{eq:congestion_state}
\end{equation}

These mechanisms allow the system to dynamically respond to 
changing network conditions in real time.

\begin{algorithm}[H]
\caption{Congestion Level Estimation}
\begin{algorithmic}[1]
\State \textbf{Input:} Queue state $Q_{current}$, $Q_{max}$, threshold $c_{th}$
\State \textbf{Output:} Congestion factor $c$, Congestion state $\mathcal{C}$

\State Compute $c \leftarrow \dfrac{Q_{current}}{Q_{max}}$

\If{$c > c_{th}$}
    \State $\mathcal{C} \leftarrow \text{Congested}$
    \State Trigger adaptive rate control
\Else
    \State $\mathcal{C} \leftarrow \text{Normal}$
    \State Maintain current transmission rates
\EndIf

\State \Return $c$

\end{algorithmic}
\end{algorithm}

%-------------------------------------------------------------
\section{Adaptive Rate Control Implementation}

The adaptive rate control module regulates the message transmission 
rate based on congestion level and message priority. Each message 
$m_i$ is assigned a priority weight $W_i$ according to its 
classified priority label $\pi_i \in \{H, M, L\}$, defined as:

\begin{equation}
W_i = 
\begin{cases}
W_H & \text{if } \pi_i = H \\
W_M & \text{if } \pi_i = M \\
W_L & \text{if } \pi_i = L
\end{cases}
\label{eq:priority_weight}
\end{equation}

\noindent where $W_H > W_M > W_L$, ensuring higher weights are 
assigned to more critical messages.

The adjusted transmission rate $R_i^{new}$ for each message is 
computed using the priority-aware rate control function:

\begin{equation}
R_i^{new} = 
\begin{cases}
R_i^{base} & \text{if } c \leq c_{th} \\
R_i^{base} \cdot \left(1 - c\left(1 - W_i\right)\right) & \text{if } c > c_{th}
\end{cases}
\label{eq:rate_control}
\end{equation}

\noindent where $R_i^{base}$ is the baseline transmission rate 
of message $m_i$. Under low congestion ($c \leq c_{th}$), messages 
are transmitted at their base rates. When congestion exceeds the 
threshold, transmission rates are reduced proportionally based on 
priority weight, ensuring high-priority messages are preserved 
while non-critical traffic is progressively throttled.

To guarantee reliable delivery of emergency messages, a minimum 
transmission rate constraint is enforced:

\begin{equation}
R_i^{new} \geq R_{min}, \quad \forall \, m_i \text{ with } \pi_i = H
\label{eq:min_rate}
\end{equation}

Priority-based message scheduling is implemented using multiple 
queues $\mathcal{Q}_H$, $\mathcal{Q}_M$, and $\mathcal{Q}_L$ 
that store messages according to their priority levels. The 
aggregate bandwidth constraint is enforced as:

\begin{equation}
\sum_{i=1}^{N} R_i^{new} \leq B_{available}
\label{eq:bandwidth_constraint}
\end{equation}

\noindent where $B_{available}$ is the total available channel 
bandwidth and $N$ is the total number of active messages.

Messages are inserted into the appropriate queue based on 
classification results. The scheduler processes messages in 
order of priority, ensuring that critical messages are 
transmitted first. Queue management operations such as 
insertion, deletion, and message selection are implemented 
using standard data structures.

The priority queue scheduling policy is expressed as:

\begin{equation}
\mathcal{Q}_{schedule} = \mathcal{Q}_H \succ \mathcal{Q}_M \succ \mathcal{Q}_L
\label{eq:priority_queue}
\end{equation}

This scheduling mechanism improves Quality of Service (QoS) by 
reducing transmission delays for critical messages \cite{rahman2022uavqos}.

\begin{algorithm}[H]
\caption{Priority-Aware Congestion Control and Scheduling}
\begin{algorithmic}[1]
\State \textbf{Input:} Labeled message stream $\mathcal{M}$ with 
       priority $\pi_i$, congestion factor $c$, threshold $c_{th}$
\State \textbf{Output:} Priority-aware updated transmission rate

\For{each message $m_i \in \mathcal{M}$}
    \State Assign weight $W_i$ based on Eq.~(\ref{eq:priority_weight})
    \State Estimate congestion factor $c$ (Eq.~(\ref{eq:congestion_factor}))
    
    \If{$c \leq c_{th}$}
        \State $R_i^{new} \leftarrow R_i^{base}$
    \Else
        \State $R_i^{new} \leftarrow R_i^{base} \cdot (1 - c(1 - W_i))$
               (Eq.~(\ref{eq:rate_control}))
    \EndIf
\EndFor

\end{algorithmic}
\end{algorithm}

This adaptive mechanism ensures that the aggregate transmission 
rate remains within the available bandwidth, thereby satisfying 
the constraint defined in Eq.~(\ref{eq:bandwidth_constraint}). 
By proactively regulating message rates, the system effectively 
prevents queue overflow and reduces packet loss during heavy 
traffic conditions \cite{khan2020mlnetwork}.

%-------------------------------------------------------------
\section{Message Transmission Simulation}

The overall Packet Delivery 
Ratio (PDR) of the system is computed as:

\begin{equation}
\text{PDR} = \frac{\text{Total Packets Delivered}}{\text{Total Packets Generated}} \times 100\%
\label{eq:pdr}
\end{equation}

The simulation environment models the following network parameters:

\begin{itemize}
\item Message generation rate
\item Network bandwidth
\item Transmission delay
\item Packet loss conditions
\end{itemize}



The system records performance parameters such as packet delivery 
success rate, average delay, and throughput. These performance 
values are later used to evaluate system efficiency and compare 
against the FIFO baseline in the results chapter.
%-------------------------------------------------------------







\section{System Integration}

All modules developed during implementation are integrated 
to form a complete communication management system.

The integrated workflow includes:

\begin{itemize}

\item Message generation

\item Feature extraction

\item Priority classification

\item Congestion monitoring

\item Rate adjustment

\item Queue scheduling

\item Message transmission

\end{itemize}

System integration ensures smooth interaction between 
individual modules and enables complete system 
functionality.

\section{Closed-Loop RL Scheduling and Aggregation Implementation}

The scheduling and aggregation extension is implemented as a two-stage 
feedback loop. One component acts as the decision layer and selects an 
action from the learned Q-table; the second component acts as the 
environment layer and simulates the effect of that action on the current 
traffic state.

The combined action has two parts:

\begin{itemize}
\item \textbf{Scheduling choice}: select which priority queue is served next
\item \textbf{Aggregation choice}: select how many packets are bundled before transmission
\end{itemize}

The implementation follows the structure below:

\begin{algorithm}[H]
\caption{Closed-Loop RL Scheduling and Aggregation Execution}
\begin{algorithmic}[1]
\State \textbf{Input:} Current network state $s_t$, learned Q-table $Q$
\State \textbf{Output:} Applied action $a_t$ and updated state $s_{t+1}$
\State Observe current queue, channel, and loss context
\State Encode the state into the learned representation
\State Select action $a_t = \arg\max_a Q(s_t, a)$
\State Decode $a_t$ into scheduling class and aggregation level
\State Apply the action to the communication simulator
\State Measure delivered bytes, attempted bytes, and loss ratio
\State Generate the next state $s_{t+1}$ from the updated traffic snapshot
\State Feed $s_{t+1}$ back to the policy for the next decision
\end{algorithmic}
\end{algorithm}

This implementation is intentionally lightweight so that the policy can 
be evaluated quickly on recorded traffic traces while still preserving the 
closed-loop control structure of the deployed system.

The policy layer is responsible only for choosing the joint action. The 
environment layer handles packet transmission, loss accounting, and the 
construction of the next state. This separation keeps the code modular and 
makes it easier to compare different trained checkpoints using the same 
simulation pipeline.

\subsection{RL Training Procedure}

The tabular Q-learning policy is trained offline on finite traffic episodes
before being reused for evaluation. Training begins with an initially empty
Q-table, and the agent explores the action space using an $\epsilon$-greedy
strategy to balance exploration and exploitation.

The training configuration used in this project is summarized below:

\begin{table}[H]
\centering
\small
\begin{tabular}{|l|c|p{6.2cm}|}
\hline
\textbf{Parameter} & \textbf{Value} & \textbf{Purpose} \\ \hline
Learning rate $\alpha$ & 0.10 & Controls how strongly each new reward update changes the Q-value \\ \hline
Discount factor $\gamma$ & 0.95 & Gives high weight to future returns and long-term queue stability \\ \hline
Exploration rate $\epsilon$ & 0.20 & Forces limited random exploration during early training \\ \hline
Number of episodes & 50 & Total offline training passes over the traffic traces \\ \hline
Episode length & Up to 500 steps & One step per scheduling decision on the replayed trace \\ \hline
Convergence behavior & 10--20 episodes & Empirical range in which the policy begins to stabilize \\ \hline
\end{tabular}
\caption{Tabular Q-Learning Training Configuration}
\label{tab:rl_training_config}
\end{table}

For each transition $(s_t,a_t,r_t,s_{t+1})$, the Q-value is updated using the
rule in Eq.~(\ref{eq:q_learning_update}). The learning rate controls how
quickly new observations override prior estimates, while the discount factor
captures the value of future returns. The exploration rate is reduced over
episodes so that the learned policy becomes more stable after the early
training phase.

In this project, one episode corresponds to one complete replay of a recorded
traffic trace. The environment returns a new state after every scheduling
decision, and the episode ends when the trace is exhausted or all queued
packets for that segment have been processed. This setup provides a clear
training definition without requiring live integration into the network
stack.

The NS-3 component is used as the traffic-generation source for the recorded
traces rather than as a fully coupled online simulator. As a result, the
implementation uses NS-3 outputs for reproducible replay while the policy
itself is trained and evaluated in the closed-loop trace-driven pipeline.

\subsection{NS-3 Dataset Generation Procedure}

The NS-3 simulator is used to generate trace files that serve as input to the
offline RL training and evaluation pipeline. The generator creates a two-node
UAV-to-ground communication setup, replays traffic at fixed time steps, and
logs one CSV row per step with queue, channel, mobility, energy, and delivery
features.

For the publication-oriented dataset set, the same generator is executed under
multiple predefined scenarios with different arrival rates, burst patterns,
mobility settings, propagation losses, and energy budgets. Each scenario is run
with its own seed so that the resulting traces remain reproducible while still
covering diverse network conditions.

\begin{algorithm}[H]
\caption{NS-3 Trace-Driven Dataset Generation}
\begin{algorithmic}[1]
\State \textbf{Input:} Scenario parameters, random seed, number of rows, step interval, output path
\State \textbf{Output:} CSV trace containing queue, link, energy, and aggregation features
\State Initialize a two-node network: UAV sender and ground receiver
\State Configure mobility, Wi-Fi channel, propagation loss, and energy models
\State Install a packet sink on the receiver and connect RX/SNR trace callbacks
\State Create the sender application and open the CSV output file
\For{each simulation step until the requested number of rows is reached}
    \State Sample new packet arrivals using a Poisson process with occasional burst arrivals
    \State Generate packet size, priority class, and deadline for each queued packet
    \State Decrease packet deadlines and drop expired packets
    \State Observe queue length, current goodput, SNR, mobility distance, and energy level
    \State Select the highest-priority packet available for transmission
    \If{queue size is large and bandwidth or SNR is below threshold}
        \State Aggregate multiple packets and transmit the combined payload
    \Else
        \State Transmit the selected packet without aggregation
    \EndIf
    \State Record packet loss ratio, link stability, AoI, and aggregation flags in the CSV row
\EndFor
\State Stop the simulator and close the trace file
\end{algorithmic}
\end{algorithm}

For the nine-scenario publication set, the outer execution loop repeats this
procedure across the scenario list and assigns a distinct seed to each run.
The helper script collects all CSV files into a timestamped output directory and
writes a manifest that records the scenario, seed, row count, and file path.

%-------------------------------------------------------------

\section{Testing and Validation}

After implementation, the system is tested under 
different traffic conditions to verify its performance using the Dataset.

Validation parameters include:

\begin{itemize}

\item Packet Loss Rate

\item Message Delivery Ratio

\item Average Transmission Delay

\item Throughput

\end{itemize}

The collected performance data is analyzed to ensure 
that the proposed system improves communication 
efficiency compared to traditional fixed-rate 
communication methods.

%-------------------------------------------------------------

\section{Summary}

This chapter presented the implementation details of 
the proposed Machine Learning-based message 
prioritization and congestion-aware rate control 
system. The implementation included dataset preparation, 
Machine Learning model training, congestion monitoring, 
adaptive rate control, priority-based scheduling, and the 
closed-loop reinforcement learning execution flow for joint 
scheduling and aggregation.

The developed system provides a complete framework for 
intelligent communication management in UAV networks. 
The performance of the implemented system is evaluated 
in the next chapter using simulation results and 
performance metrics.